\documentclass[../main.tex]{subfiles}
\begin{document}
%==========================================TIÊU ĐỀ TẬP========================================
\begin{table}[h]
  \begin{adjustwidth}{-1cm}{}
    \begin{tabular}{>{\centering\arraybackslash}p{6cm}|>{\raggedright\arraybackslash}p{12.5cm}}
      \multirow{5}{6cm}
      % Cột bên trái: ÉPISODE và số tập
      {\\[-40pt]
        \flaregothic\fontsize{20pt}{20pt}\selectfont \centering
        \color{eptype}ÉPISODE\color{black}\\
        \burbank\fontsize{72pt}{72pt}\selectfont
        \color{epnum}92
        \\[18pt]
        \flaregothic\fontsize{14pt}{14pt}\selectfont
        \color{epattr}BIENVENUE, \uline{\color{Shishiron} Botan} !
        \flaregothic\fontsize{18pt}{18pt}\selectfont
        \color{epattr}ÉPISODE PERDU RESTAURE\\
        \color{black}
      }
       &
      % Dòng thứ nhất: tên tập tiếng Nhật
      {\fotskip\fontsize{18pt}{18pt}\selectfont \ruby{三}{さん}\ruby{輪}{りん}\ruby{車}{しゃ}に\ruby{乗}{の}る\ruby{大}{おお}\ruby{人}{とな}たち
      } \\[6pt]
      % Dòng thứ hai: tên tập tiếng Anh
       & {\cafeteria\fontsize{18pt}{18pt}\selectfont The Tricycle}                                                         \\[4pt]
      % Dòng thứ ba: tên tập tiếng Việt
       & {\vnmsans\fontsize{18pt}{18pt}\selectfont Chuyện về xe ba bánh}                                                   \\[8pt]
      % Dòng thứ tư: Nhân vật xuất hiện
       & {\brian{\fontsize{14pt}{14pt}\selectfont Track used: Jurassic Marsh - Ultimate Battle}
         }
      \\[8pt]
      % Dòng cuối cùng: Thời gian diễn ra của tập (thường sẽ là ngày phát hành)
       & {\continuum\fontsize{16pt}{16pt}\selectfont Ngày phát hành: 07/02/2021}                                           \\[18pt]
    \end{tabular}
  \end{adjustwidth}
\end{table}
%=========================================TRUYỆN CHÍNH========================================
\color{black}

\pov{Hikawa Kuon}

Thành viên thứ ba, và cũng là thành viên cuối cùng (tính từ thời điểm này) sẽ đặt chân tới văn phòng \hll\ trong ngày hôm nay.

Một nàng sư tử trắng với phong thái điềm tĩnh, mang trong mình khí chất vương giả nhưng lại chẳng hề phô trương. Nếu Lamy là một nàng Yêu tinh tuyết dịu dàng, thì thành viên mới này lại là một chiến binh thầm lặng, luôn bình thản đối mặt với mọi thử thách như thể chẳng có gì có thể làm em ấy nao núng.

Em ấy sở hữu một bản lĩnh đáng kinh ngạc: dù thường tỏ ra lười biếng và thích tận hưởng cuộc sống, nhưng thực chất lại là một game thủ cừ khôi với phản xạ nhanh như chớp trong lĩnh vực bắn súng góc nhìn thứ nhất - lĩnh vực mà tôi dở nhất. Sự điềm tĩnh, khả năng phân tích tình huống và óc phán đoán sắc bén khiến em ấy trở thành một người đáng tin cậy trong bất kỳ cuộc chơi nào.

Ngoài ra, tính cách của em cũng rất độc đáo: thoải mái, dễ gần khi ở bên những người thân thiết như Lamy-chan hay Polka-chan, nhưng lại có phần khó nắm bắt với những ai chưa hiểu rõ em ấy. Em ấy không thích những thứ ồn ào hay phô trương, mà chỉ đơn giản tận hưởng những điều mình yêu thích theo cách riêng của mình. Một chiến binh mạnh mẽ nhưng chẳng bao giờ khoe khoang về sức mạnh của bản thân.

Bên cạnh đó, em ấy còn là một chuyên gia\dots\ ngủ nướng và sống một cuộc sống thư giãn.

Và dù là một con sư tử, em không hề tìm kiếm vinh quang trên chiến trường. Thay vì vậy, em muốn một ngày nào đó có thể đứng trên một sân khấu lớn, nơi mọi người không chỉ nhìn thấy sức mạnh mà còn cảm nhận được niềm vui và sự bình yên mà em mang lại.

Một sư tử với giấc mơ phiêu du.

\dots Hoặc có lẽ, một ``ẩn sĩ'' giấu mình dưới lớp vỏ của một kẻ lười biếng.

\chrint

Shishiro Botan (\ruby{獅}{し}\ruby{白}{しろ}ぼたん, \ruby{Si}{Sư}-\ruby{si-rô}{Bạch}\ Bô-tan) là thành viên tiếp theo và cuối cùng (ít nhất là cho tới hiện tại) của thế hệ thứ Năm. Một nữ ``chúa tể sơn lâm'', chỉ khác ở chỗ thay vì ``nữ hoàng'' ở ngoài xa van tự nhiên thì em ấy là ``nữ hoàng'' của những\dots\ trò chơi bắn súng góc nhìn thứ nhất.

\begin{wrapfigure}[27]{r}{7.5cm}
  \includegraphics[width=7.5cm]{../../../../Image/Book?/botan_model.png}
\end{wrapfigure}

Botan-chan là một người có thể nói là ``ngoại hình và tính cách hoàn toàn không ăn khớp với nhau.'' Trái ngược với lý lịch lười biếng và vô tư của mình, em ấy có tinh thần trách nhiệm cao và thích môi trường làm việc có cấu trúc. Vẻ ngoài lạnh lùng, trưởng thành của Botan-chan trái ngược với bản tính vui vẻ và dễ tính. Em có một tông giọng khá cao, khác xa hoàn toàn với tưởng tượng của tôi khi đọc sơ yếu lý lịch lần đầu.

Em ấy là người rất thoải mái, kể cả việc khi em trò chuyện với chúng tôi hay khi lên sóng, và đây là một trong những tiêu chí quan trọng để làm cho người xem có thể phấn chấn nhất. Mọi người trong thế hệ thứ Năm đều coi em ấy là thủ lĩnh của nhóm, quan sát và giúp đỡ những người khác với tư cách là game thủ giàu kinh nghiệm nhất.

Mặc dù Botan-chan là một người có tính nghiêm túc và giàu động lực, em ấy lại\dots\ thích trở thafanh một người ngốc nghếch, thể hiện xu hướng thích chơi đùa nhiều hơn khi ở cạnh những người thuộc tiền bối, trong đó có cả tôi.

Em ấy rất quan tâm đến những thành viên \hll\ đồng thế hệ của mình, nhất là Lamy-chan, một cặp đôi bài trùng không thể tách ròi, và đôi khi đóng vai trò là người lãnh đạo nhóm. Em ấy - giống như Mio-chan hay Fubu-chan, cũng\dots\ tỏ vẻ muốn ăn thịt Watame-chan, dù tôi còn chẳng biết em ấy đang đùa hay đang có ý định thật. Em ấy cũng là một người nghịch ngợm, thể hiện việc hay kéo mọi người vào chơi trò kinh dị, hay cười nhiều, nhất là khi em ấy gặp khó khăn.

Tiếng cười độc nhất vô nhị của em ấy là một đặc điểm được người hâm mộ của cô đánh giá cao, ngoài ra em ấy còn có những tiếng động nghe khá vui tai, ví dụ: ``Poi.''

Em ấy có đôi tai sư tử đeo khuyên, đuôi sư tử, đôi mắt xám, hàm răng nanh và mái tóc dài màu xám bù xù được buộc hai bên với tóc dài hai bên, tóc mái tạo hình chữ A và phần mái chia đôi dài đến giữa hai mắt.

Botan-chan thường diện một chiếc áo sơ mi bó sát không tay màu đen, hở khe ngực bằng một tấm màn xuyên thấu, đeo một chiếc vòng cổ đen, mặc một chiếc áo khoác lông màu đen mặc trễ vai, một chiếc váy ngắn có khóa kéo màu xám. Em có đi một chiếc quần tất đen rách ở đùi trái, một chiếc dây đeo đùi hình chữ O và đi bốt cao gót màu đen.

\chrint

Với bản sơ yếu lý lịch đồ sộ này, không biết cô nàng Sư tử này sẽ có trò gì để đối phó với tôi đây\dots

\il

Chủ Nhật đầu tiên của tháng Hai Dương lịch. Với tôi, Tết Nguyên Đán năm Tân Sửu đã khá cận kề - hôm nay đã là 26 tháng Chạp Âm lịch. Không khí những ngày cuối năm Âm lịch lúc nào cũng khiến tôi có cảm giác bồi hồi, dù ở Nhật Bản không tổ chức Tết theo lịch Âm.

Hôm nay là ngày làm việc cuối cùng của tôi trong năm Canh Tý Âm lịch, và cũng là ngày mà Botan-chan lần đầu đặt chân đến văn phòng công ty, trước khi tôi bắt đầu chính thức nghỉ Tết.

Nàng Sư tử trắng bước vào với dáng vẻ uể oải, tay dụi mắt, miệng thì ngáp dài một hơi rõ to. Tôi nhìn em ấy rồi thở dài: ``Nhìn em ngáp thế này trông anh cũng muốn đi ngủ quá\dots''

Tôi vươn vai một chút rồi chống tay lên bàn. ``Nhưng mà làm sao mà ngủ nổi khi mà trời vẫn còn đang sáng chưng thế này.''

``Anh nói thế nào chứ'' Botan-chan vừa nói vừa vươn vai lên cao, ``hôm nay là một ngày đẹp đấy. Chắc em cũng nên ra ngoài đi dạo một lúc nhỉ.''

\img{7-10a}

Trời đất, hiếm khi một người lười chảy thây như em ấy lại muốn được ra ngoài đấy. Nhưng ngay sau đó, tôi đổ giọt mồ hôi trên đầu, tỏ vẻ hơi ái ngại: ``Anh không nghĩ vậy đâu.''

Em ấy nghiêng đầu, vẻ ngạc nhiên: ``Ủa, sao lại thế?''

``Anh xem dự báo thời tiết sáng nay rồi. Tầm tí nữa là mưa bây giờ đấy.''

Botan-chan chớp mắt, sau đó nhún vai, ngó ra ngoài cửa sổ: ``Làm sao mà mưa được nhỉ. Trời đang đẹp mà.''

Tôi chẹp miệng, lắc đầu: ``Để xem.''

Và như thể trời cũng muốn chứng minh lời tôi nói, ngay sau khi tôi vừa dứt câu, những đám mây xám ùn ùn kéo tới, rồi mưa bắt đầu trút xuống ào ào. Sấm chớp vang lên đùng đoàng, làm rung cả cửa kính. Botan-chan vừa bị tụt hứng, vừa lườm tôi với ánh mắt khó chịu.

\img{7-10b}

Tôi nhún vai, tặc lưỡi: ``Bảo rồi mà.'' Cơ mà, cơn mưa đến một cách bất chợt như vậy làm tôi cũng\dots\ không bất ngờ lắm, vì\dots\ chắc các bạn cũng biết rồi đấy.

Cả hai chúng tôi chỉ còn biết đứng bên cửa sổ, nhìn cơn mưa xối xả ngoài trời trong sự bực bội và chán chường.

\ct

Đã ba mươi phút trôi qua. Trời vẫn mưa xối xả bên ngoài, chẳng có dấu hiệu gì là sẽ tạnh sớm. Thậm chí, có vẻ như trời càng lúc càng mưa to hơn.

Mùa đông mà mưa rào như mùa hè thế này quả đúng là chuyện lạ.

Tôi ngồi trên ghế sô-pha, tay cầm điện thoại bấm bấm cho đỡ chán. Ở phía cửa sổ, nàng Sư tử thì lại vác cây súng bắn tỉa lên vai, với phần mũi súng chúi xuống dưới, đứng tạo dáng đầy phong cách.

\img{7-10c1}

Tôi thì không sợ khẩu súng đó, có hai lý do. Một, tôi biết rõ khẩu súng ấy thực chất chỉ là mô hình. Hai, đặc trưng của người từ quê tôi - mấy thứ vũ khí như thế này đối với tôi mà nói - nếu là hàng thật thì cũng chỉ như một món đồ chơi.

``Nhìn thế chứ trông em cũng ngầu phết đấy,'' Tôi gật gù, vỗ tay nhẹ như thể đang đánh giá một tác phẩm nghệ thuật. Nhưng Botan-chan vẫn giữ nguyên nét mặt lạnh tanh, chẳng có vẻ gì là phản ứng lại cả.

Thế nhưng, có ai đó vừa mới bước vào phòng lại không nghĩ như chúng tôi.

Núp sau cái kệ tủ là Towa-chan, người đang run lên bần bật, mặt thì tái mét, toát mồ hôi hột. Nàng Sư tử trắng nhìn thấy, nghiêng đầu khó hiểu: ``Senpai?''

\img{7-10c2}

Nàng Ác ma nuốt nước bọt cái ực, run rẩy với khẩu súng mô hình đằng sau y như thật giắt sau lưng của nàng Sư tử, cùng với bộ dạng hầm hố của em ấy.

Tôi liếc qua, thở dài: ``Có gì đâu mà phải sợ. Cứ ra kia nói chuyện bình thường, em ấy không cắn đâu.''

Botan-chan chớp mắt, vẫn không hiểu chuyện gì đang xảy ra. ``Sao chị lại né tránh em thế?''

Towa-chan vẫn đứng im như tượng, mắt dán chặt vào khẩu súng mô hình trên vai của đàn em, lắp bắp: ``Em\dots\ em đang xách cái súng to bỏ bố kia kìa! Kuon-senpai, anh không thấy sợ à!?''

Tôi nhún vai, đáp tỉnh bơ: ``Anh thấy bình thường mà.''

``Thôi nào chị,'' Botan-chan cười khổ, ``chị đang hiểu nhầm rồi.''

``Ít nhất thì cũng phải để em ấy giải thích chứ,'' tôi thở hắt ra.

Nhưng chưa kịp để chúng tôi giải thích, nàng Ác ma đã ôm lấy ngực mình, giọng đầy kịch tính: ``Nhưng mà, trái tim chị chỉ toàn sự bi quan thôi\dots''

Tôi chống nạnh, lắc đầu đầy ngán ngẩm: ``Ác ma Towa-sama gì mà nhát thế\dots''

Nàng Sư tử đứng bên cạnh, hơi ngạc nhiên: ``Hả?''

``Không phải em nhát đâu!'' nàng Ác ma lập tức phản đối. ``Em luôn chuẩn bị cho tình huống xấu nhất mà!''

Tôi chẹp miệng: ``Nhưng ít nhất em cũng phải nghĩ tích cực một chút chứ!''

Đến đây, có vẻ như sự ``đồ sộ'' và ``lạnh lùng'' của Sư tử khiến nàng Ác ma không còn chịu nổi nữa. Towa-chan bất thình lình xoay người bỏ chạy ra khỏi phòng, không buồn đoái hoài gì tới những lời chúng tôi gọi với theo.

\img{7-10d}

\begin{dialogue}
  \speak{Botan} Đ-Đợi đã nào, senpai!
  \speak{Kuon} Để cho em ấy cơ hội để giải thích chứ!
\end{dialogue}

*RẦM!*

Tiếng cửa mở ra rồi đóng lại.

Botan-chan chớp mắt nhìn tôi, rồi cúi đầu xuống, giọng có chút buồn bã: ``Chỉ là em\dots\ giữ nó như vậy cho đến khi keo khô thôi mà\dots''

Khẩu súng bắn tỉa mô hình bằng nhựa mà em ấy đang cầm chính là do tự tay em ấy lắp ráp. Em ấy đã trát keo lên một số bộ phận, và chính tôi là người bảo em ấy nên giữ như vậy để keo nhanh khô cứng hơn.

Tôi khoanh tay, tặc lưỡi. Chắc là Towa-chan nhát quá thôi. ``Thế mà cũng tự xưng mình là Ác ma được'' - tôi thật sự muốn buột miệng mắng em ấy một câu, nhưng thôi, tôi sẽ để dành nó vào một dịp khác vậy.

\ct

Một lúc sau, tiếng cửa phòng lại vang lên. Sư tử trắng tưởng rằng tiền bối của cô quay trở lại nên đã thốt lên đầy hy vọng: ``Towa-senpai!''

Nhưng không, người xuất hiện không phải là nàng Ác ma.

Thay vào đó, bước vào phòng là một cô bé\dots\ không, phải nói chính xác hơn là một Công chúa đang đeo kính râm, cưỡi trên một chiếc xe đạp ba bánh màu đỏ lấp lánh.

\img{7-10e}

Tôi trợn mắt, vô thức thốt lên: ``Oắt đờ phắc\dots? Luna-chan, bộ dạng này là sao đây?''

Botan-chan cũng chớp mắt đầy khó hiểu: ``Ờm\dots\ chị có ổn không vậy?''

Luna-chan ngẩng cao đầu, chống nạnh đầy uy nghi rồi bật ra một câu còn khó hiểu hơn: ``Em có ván đề zì mà em mún được zải quýt hơm-nora?''

Hai người chúng tôi nhìn nhau, chẳng ai hiểu nổi cái tình huống này là gì nữa. Trông em ấy như thể là một đứa trẻ mầm non đang tập làm xã hội đen vậy.

Nàng Sư tử hắng giọng, hơi nghiêng đầu. ``Thì\dots\ vừa nãy em có đấy.''

``Vậy thì em có hỏy các chụy ở đây thôi, Botan-chan,'' nàng ``Em bé'' phẩy tay. ``Chụy có thể giúp em được. Mà chị tưởng Kuonsa-senpai giúp được em mà?''

Tôi khoanh tay, lắc đầu. ``Chả biết một em bé có thể giúp được gì cho một con Sư tử nữa\dots''

``Đừn có mừ khinh thừn em chứ,'' Luna-chan đáp lại, lườm tôi với ánh mắt bất bình.

Botan-chan thở dài, nhìn về phía đàn chị của mình: ``Em cũng thử rồi đấy, nhưng mà\dots''

Tôi tiếp lời, nhún vai đầy bất lực: ``Em ấy nhát như cáy, còn chưa kịp để bọn anh giải thích nữa.''

Nàng Công chúa xứ Kẹo ngồi vắt chân trên xe ba bánh, tay gõ gõ vào tay lái. Trông em ấy có vẻ vẫn chưa hiểu chuyện gì đang xảy ra.

Tôi với Botan-chan nhìn nhau, rồi cùng thở dài. Thôi thì giải thích từ đầu vậy.

``Chuyện là thế này\dots''

\ct

Botan-chan buông một tiếng thở dài, kết thúc câu chuyện hồi tưởng của hai chúng tôi: ``\dots và sau đó, em nghĩ Towa-senpai đã hiểu nhầm ý của em, nên em muốn nhờ chị giải thích với em ấy chuyện đó. ''

Nghe xong đầu đuôi câu chuyện, Luna-chan không trả lời ngay, mà chỉ lẳng lặng bước tới cửa sổ, nhìn ra ngoài trời mưa rồi cất giọng đầy bí ẩn: ``Nhìn ra ngoày cử xổ đi-nora.''

Tôi và nàng Sư tử nhìn nhau, rồi cũng làm theo.

Mưa vẫn đang trút xuống không ngừng. Nhưng ở ngoài kia\dots\ có một bóng người nhỏ bé đang di chuyển trên con đường ướt sũng.

Hoặc đúng hơn, một bóng người đang đạp xe ba bánh ngoài kia.

\img{7-10f}

Tôi nheo mắt nhìn kỹ hơn, rồi buột miệng: ``Kia chả phải là Marine-chan à?'' Trông em ấy như thể đang bị mê sảng vậy: đôi mắt thì xoay như chong chóng, miệng thì lẩm bẩm cái gì đó, chân thì đạp đầy lảo đảo, nhìn như có thể gục bất kỳ lúc nào.

Nàng Sư tử cũng nhíu mày khó hiểu. ``Chị ấy đang làm gì ở ngoài này vậy?''

Tiền bối (về lý thuyết) của em ấy nhún vai, điệu bộ như thể đang giảng giải một chân lý hiển nhiên: ``Em có bít chụy ấy ngĩ thế nào khi mà một cứp biển lại đi đạt se ba bánh hơm?''

Em ấy chớp mắt mấy lần, rồi thành thật đáp: ``Chịu chết luôn.''

Tôi lắc đầu, nhìn cảnh tượng ngoài kia mà chẳng biết nên cười hay nên khóc. ``Tự dưng dở hơi, trời mưa lại đi đạp xe\dots\ Lại còn xe ba bánh nữa. Bộ em nghĩ em ấy là nhi đồng thối mũi thích tắm mưa chắc?''

Luna-chan cười khúc khích, rồi hất cằm đầy tự tin: ``Chụy í đang ngĩ zằng mình đang ở chên tàu cứp biển đấy-nanora.''

Sư tử trắng xanh mặt với vẻ ngỡ ngàng còn tôi thì chỉ biết há hốc mồm rồi phản bác: ``Nó chả có lý gì cả!''

Nhưng nàng Công chúa xứ Kẹo không hề nao núng, chỉ chậm rãi đặt tay lên vai của nàng Sư tử và nghiêm túc giảng giải ``triết lú'': ``Đó, hỉu nhầm về một cái gì đó là như vậy á-nora.''

Botan-chan thì nhìn Luna-chan, rồi nhìn tôi, cuối cùng buông một tiếng thở dài bất lực: ``Giờ em muốn về nhà và chơi game quá\dots''

Tôi đổ mồ hôi hột. ``Cái đó gọi là `phê đá' mới đúng đó, Luna-chan ạ.''

``Thì nó cũng là hỉu nhầm thôi na.''

Tôi khoanh tay, lắc đầu: ``Nó khác hoàn toàn đấy. Đúng là suy nghĩ đơn sơ của em bé có khác.''

Em ấy chỉ nhún vai đáp lại: ``Tuỳ anh ngĩ thế nào thui. Mà, Botan-chan\dots''

Gọi nàng Sư tử xong, em ấy đứng dậy khỏi chiếc xe ba bánh, rồi nghiêm túc nói: ``Nếu em nái chiếc xe này, em sẽ hỉu zõ hơn đó-nora.''

Tôi nheo mắt đầy nghi ngờ: ``Cái này thì hiểu rõ hơn được gì?''

``Cứ nái đi thì sẽ bít na!'' Luna-chan vẫn khăng khăng.

Botan-chan lùi lại một chút, vẫy tay từ chối: ``Không, chị ơi\dots\ Nói thật với chị là em cũng không muốn lái đâu.''

Nhưng đáp lại chỉ là một cái chỉ tay dứt khoát ra ngoài cửa sổ của ``Em bé'':

``Em lại nhìn ra ngoài cửa xổ cho chị-nora.''

Chúng tôi nhìn theo. Marine-chan vẫn đang miệt mài đạp xe giữa cơn mưa, và vẫn chưa có dấu hiệu gì là tỉnh táo.

``Thế nào?'' nàng tóc hồng khoanh tay, nhướng mày hỏi.

Tôi chỉ biết thở dài đầu hàng trước suy nghĩ ``ngây thơ vô số tội'' của em ấy.

Nàng Sư tử cũng im lặng một lúc, rồi cuối cùng chịu nhún nhường: ``\dots Thôi được, em sẽ lái.''

\ct

Thế là nàng Sư tử - biểu tượng của sự dũng mãnh - lại ngồi trên một chiếc xe đạp ba bánh, thứ mà bình thường chỉ có trẻ con mới dùng. Sự đối lập đến mức ngược đời, nhưng cũng\dots\ buồn cười một cách kỳ lạ.

Cơ mà nghĩ lại, có khi sự ngược đời cũng mang lại nét đáng yêu riêng.

\img{7-10g}

Botan-chan khoanh tay, nhìn xuống chiếc xe nhỏ bé mà mình đang ngồi lên, rồi hỏi: ``Với cái này thì em sẽ giải quyết kiểu gì?''

Tôi chỉ nhún vai, đáp tỉnh bơ: ``Anh chịu. Cái suy nghĩ đơn giản của em ấy thì đến cụ anh còn chả hiểu nổi, huống gì đến em.''

Trong lúc chúng tôi đang cố lý giải cái logic chối tai mà Luna-chan đã bày ra, bỗng nhiên, cánh cửa lại bật mở.

Người vừa bước vào chính là nàng Ác ma - người mà nàng Sư tử đang muốn thanh minh với. Towa-chan vội vàng cúi đầu xin lỗi: ``Xin lỗi em nha, Shishiron! Nghĩ lại thì, nó không thể nào là súng thật được\dots''

Botan-chan ngước lên, vui mừng gọi: ``Towa-senpai!''

Tôi khoanh tay, hất cằm: ``Thấy chưa? Anh định nói với em như thế đấy, vậy mà tự dưng chạy mất dép luôn.''

Towa-chan cười trừ, gãi đầu: ``Em cũng đã hoảng quá rồi-''

Nhưng chưa kịp nói hết câu, Towa-sama bỗng khựng lại.

Mắt em ấy mở to, nhìn chằm chằm vào đàn em của mình.

Không ai nói gì.

Không gian bỗng dưng trở nên yên tĩnh lạ thường.

Chỉ có tiếng mưa vẫn rơi rả rích bên ngoài cửa sổ.

\dots Và sau vài giây im lặng, nàng Ác ma bất ngờ lấy điện thoại ra.

*TÁCH!*

\img{7-10h}

Âm thanh quen thuộc của máy ảnh điện thoại vang lên.

Nàng Sư tử bàng hoàng nhận ra sự thật: cô vừa bị chụp ảnh trong tư thế ``ngồi xe ba bánh,'' một tư thế tượng trưng cho trẻ con, trái ngược hoàn toàn với sự dũng mãnh của em ấy.

Botan-chan hấp tấp vươn tay ra, hốt hoảng: ``Đợi đã, Senpai!''

Nhưng nàng Ác ma đã lùi về sau một bước, ung dung lướt trên màn hình điện thoại. ``Mình phải tweet cái này mới được.''

Tôi đứng một bên, khoanh tay quan sát mà không giấu nổi nụ cười thích thú: ``Nhìn thế trông cũng đáng yêu đấy.''

Còn ``nhân vật chính'' của bức ảnh thì hoảng loạn hơn bao giờ hết: ``\textbf{Đừng! Ý của em không phải như vậy đâu!}''

\dots Tuyệt lắm, Công chúa ạ. Không những không giải quyết được chuyện, mà em còn tạo thêm rắc rối nữa.

%==========================================TRUYỆN PHỤ=========================================
\omk

Và thế là, trong một buổi chiều mưa gió, tôi phải chứng kiến cảnh tượng có lẽ chỉ xuất hiện trong phim hài Tết: một con Sư tử cầm xe ba bánh đuổi theo một Ác ma đang chạy bán sống bán chết. Vâng, là \emph{cầm}, không phải \emph{đạp}.

\gif{7-10i}{1}{84}

Towa-chan hét lên, giọng đầy hoảng hốt: ``\textbf{Đừng có đuổi theo chị chứ!}''

Còn Botan-chan thì vừa chạy vừa thốt lên trong tuyệt vọng: ``\textbf{Senpai! Xóa cái bức ảnh đó đi mà\dots!}''

Với sức mạnh của một chiến binh FPS thực thụ, nàng ta vác luôn cả chiếc xe ba bánh mà lẽ ra chỉ dùng để ngồi, giờ lại thành một vũ khí đáng sợ vung qua vung lại. Tôi không biết em ấy định làm gì, nhưng nhìn thế này thì có vẻ Towa-chan đang thấy mạng sống của mình bị đe dọa.

Tôi thở dài, quay sang nhìn thủ phạm gián tiếp gây ra vụ này: ``Thấy chưa, Luna-chan? Tự dưng lại tạo thêm rắc rối cho hai người ấy rồi đấy.''

Nhưng Luna-chan lại tỉnh bơ đáp: ``Đó cũm có thể gọi là xự hỉu nhầm na.''

Tôi khựng lại, ngẫm nghĩ một lúc lâu, rồi khoanh tay, nhíu mày: ``Cái đó anh phải công nhận\dots\ Mà cái này là do em làm đấy.'' Em ấy chỉ đáp lại bằng một cái nhún vai, tỏ vẻ không can tâm gì.

Mặc kệ cuộc trò chuyện của chúng tôi, hai người kia vẫn tiếp tục màn rượt đuổi kịch liệt khắp phòng. Botan-chan, với một tay cầm xe đạp ba bánh, không ngừng vung vẩy trong khi đuổi theo đối phương.

``\textbf{Đứng lại đó! Mọi chuyện không phải như chị nghĩ đâu!}'' em ấy thét lên, cố gắng thanh minh. Ngược lại, đàn chị của em ấy thì vừa ôm đầu, vừa sợ hãi đáp lại: ``\textbf{Tha cho chị đi mà!}''

Tôi bắt đầu cảm thấy lo lắng thật sự. Nếu còn tiếp tục thế này, rất có thể phòng nghỉ của \hll\ sẽ sớm trở thành hiện trường vụ án mất.

Tôi quay sang hỏi người đang đứng nhìn mọi chuyện với vẻ thích thú: ``Em không định ngăn cản hai người đó à?''

Nhưng Luna-chan lại đáp một câu hết sức ``triết lú'' với cặp kính râm vẫn đang đeo, che giấu đi ánh mắt thích thú: ``Quả nhin em í đúng là kẻ xăn mồi na.''

Tôi cạn lời. Nhìn cái tướng với mặt thế kia là không có cách gì để giúp rồi.

Không thể để mọi chuyện tiếp diễn thêm, tôi đành bất đắc dĩ ra tay.

Chỉ trong khoảnh khắc, tôi lao tới chặn trước mặt nàng Sư tử, giơ tay giữ lấy chiếc xe ba bánh trước khi nó có thể vung xuống ai đó. ``Dừng lại nào, đừng có gây thêm rắc rối nữa!''

Và đúng như tôi lo sợ, nàng Sư tử - vì quá tập trung vào cuộc truy đuổi và đang cố gắng bào chữa bản thân - suýt nữa đã ``quất'' thẳng cái xe vào tôi.

Cái này đúng là quá sức chịu đựng rồi\dots

\pov{Shishiro Botan}

Tôi vẫn còn đang hăng máu, tay siết chặt lấy chiếc xe ba bánh, cố gắng phân bua với Towa-senpai, thì đột nhiên Kuon-senpai chạy cắt ngang trước mặt tôi, giơ tay ra chắn. ``Dừng lại nào, đừng có gây thêm rắc rối nữa!''

Tôi khựng lại ngay lập tức, suýt chút nữa đã ``quất'' thẳng cái xe vào người anh ấy. Tim tôi đập thình thịch vì bất ngờ.

Tôi không giận chị ấy đâu, nhưng việc Towa-senpai cứ chạy như thế thì làm sao tôi có thể bắt chuyện được cơ chứ!

``Em đã bảo rồi, mọi chuyện không như chị nghĩ đâu!'' tôi quay sang Senpai đang nấp đằng sau anh Tổng, giọng đầy oan ức.

Nhưng người chị tiền bối của tôi lại lùi ra xa một chút, nhìn tôi đầy cảnh giác: ``Sao lại cầm chiếc xe đánh chị thế!?''

\dots Đúng thật. Có lẽ bởi vì tôi hơi quá chăng?

Kuon-senpai thở dài, xoa trán: ``Hai người cứ bình tĩnh, cứ để Botan-chan giải thích xem nào.''

Phải mất một lúc, tôi mới có thể hít thở sâu và trấn tĩnh lại. Nhìn sang Towa-senpai, tôi thấy chị ấy cũng bắt đầu dịu xuống. Tôi thả lỏng tay, đặt chiếc xe ba bánh xuống đất và bắt đầu kể lại đầu đuôi câu chuyện, từ chuyện hiểu lầm về khẩu súng mô hình cho đến màn ``đưa xe ba bánh vào triết lý sống'' của Luna-senpai.

Tôi liếc nhìn Luna-senpai và Kuon-senpai đứng quan sát mọi chuyện. ``Chị đại'' đeo kính râm thì trông rõ ràng là đang rất hứng thú với cuộc nói chuyện của hai chúng tôi, còn anh Tổng thì chỉ ngán ngẩm.

Khi tôi kể xong, anh ấy nhìn về phía người bên cạnh và nói với sự quở trách: ``Em cũng có phần trong đó đấy. Tí nữa ra xin lỗi `lính mới' đi.''

Luna-senpai gật đầu, đáp tỉnh bơ: ``Rùi\~{}'' Rõ ràng là chị ấy chẳng có chút gì là hối lỗi cả. Mà nghĩ lại, chị ấy có phải là bề trên của tôi không?

Tôi liếc nhìn Kuon-senpai, thấy anh ấy lại thở dài một hơi, nhún vai như thể còn một đống việc phải giải quyết: ``Tí nữa anh phải giải quyết chuyện của Marine-chan nữa đấy\dots''

Nghe đến đây, tôi mới sực nhớ ra cái cảnh Marine-senpai đạp xe ngoài trời mưa. Chắc bây giờ chị ấy cũng mệt lắm rồi.

Cuối cùng thì mọi chuyện cũng sáng tỏ. Towa-senpai đồng ý xóa bức ảnh chị ấy vừa chụp, nhưng mà\dots\ tôi có cảm giác hơi muộn mất rồi. Vì tôi đã lướt Twitter một chút và thấy chị ấy đã đăng ảnh tôi ngồi trên xe ba bánh lên.

Bình luận bên dưới thì\dots\ đúng như tôi dự đoán, toàn là những câu trêu chọc, dù không có ý xúc phạm hay giễu cợt gì:

\begin{itemize}
  \item ``Nhìn mặt Sư tử kìa, cưng ghê!''
  \item ``Shishiron, sao lại quay về tuổi thơ thế này?''
  \item ``Trông như một con mèo to xác bị đặt lên xe ba bánh ấy, haha.''
\end{itemize}

Tôi chỉ biết đơ người, mặt nóng bừng bừng. Giờ mà có cái hố nào trước mặt, tôi sẵn sàng nhảy xuống luôn\dots

Mà khoan đã, cái xe ba bánh này\dots\ ai là người khởi xướng vụ này ấy nhỉ?

Tôi liếc sang Luna-senpai. Không lẽ\dots\ là do chị ấy tự biên tự diễn?

\dots Ờ, thôi bỏ đi. Nghĩ nhiều lại mệt. Với cả, những người xem được bức ảnh đó cũng không lăng mạ tôi hay gì, nên chắc cũng tặc lưỡi cho qua được.

\ct

Sau đó, tôi nghe nói Kuon-senpai đã giúp Marine-senpai tỉnh lại bằng cách nào đó. Nhưng cuối cùng, chị ấy lại bị đổ bệnh vì dầm mưa quá lâu.

Hậu quả là gì ư? Chị ấy đã phải nghỉ stream mấy ngày liền vì cảm lạnh.

Chắc chắn đây là lần đầu tiên trong lịch sử, một Thuyền trưởng Hải tặc bị ``đánh chìm'' bởi một chiếc xe ba bánh. Hoặc đúng hơn, là chiêu bài của Luna-senpai.

\il

Sau tất cả những chuyện vừa diễn ra, tôi chỉ có thể thở dài và tự nhủ rằng hôm nay đúng là một ngày không bình thường chút nào.

Tôi vốn nghĩ mình chỉ cần lên văn phòng một lát để có thể tận hưởng một ngày nắng đẹp. Nhưng rồi\dots\ tôi lại phải giải thích chút chuyện với Towa-senpai, để rồi bị cuốn vào hết sự việc này đến sự việc khác. Một khẩu súng mô hình, một chiếc xe ba bánh, một thuyền trưởng cướp biển lạc lối trong cơn mưa\dots\ Đúng là không thể đoán trước được điều gì khi ở cùng mọi người trong công ty này.

Dù vậy, tôi cũng phải công nhận một điều: nếu không có Kuon-senpai đứng ra điều hành mọi chuyện, chắc tôi và Towa-senpai vẫn còn cãi nhau, Luna-senpai vẫn sẽ tiếp tục ``giảng giải triết lý'' mà tôi không thể hiểu nổi, và Marine-senpai có khi vẫn còn đạp xe ngoài trời mưa đến kiệt sức mất.

Anh ấy đúng là kiểu người mà dù có phàn nàn thế nào thì vẫn luôn làm tốt vai trò của mình. Dù bận trăm công nghìn việc, lúc nào Senpai cũng có thể lo liệu ổn thỏa mọi thứ. Và hơn hết, anh ấy hiểu rõ từng người trong chúng tôi, biết ai sẽ phản ứng ra sao, ai sẽ cần lời khuyên hay một cú hãm phanh kịp thời.

Tôi nhìn Kuon-senpai đang nhắn tin cho ai đó trên điện thoại, có vẻ là để giải quyết nốt chuyện của Marine-senpai. Nghĩ đến đây, tôi chợt nhớ ra một chuyện quan trọng.

``À mà, Senpai này?'' tôi lên tiếng, làm anh ấy ngước nhìn lên tôi.

``Sao?''

``\textit{NePoLaBo} bọn em vốn không phải chỉ có ba người.''

Anh ấy im lặng một chút rồi gật đầu: ``Anh biết chứ, thấy bốn âm tiết là đủ hiểu rồi. Vẫn còn một người nữa, đúng không?''

``Cậu ấy bảo rằng sẽ lên văn phòng vào mùa Xuân tới,'' tôi nói tiếp, ``bây giờ buổi diễn trực tiếp online sắp tới rồi, cộng với việc cô ấy cũng đang bận\dots''

Dĩ nhiên, đó là lý do mà cậu ấy đưa ra, nhưng tôi thì biết rõ sự thật hơn.

Vì cơn cảm lạnh kéo dài cả tháng nay, cậu ấy không thể đến được. Và lý do khiến cô ấy đổ bệnh\dots\ chắc cậu ấy cũng đoán được phần nào rồi nhỉ? Nhưng thôi, tôi sẽ không nói gì thêm.

Tôi chỉ hy vọng rằng, khi mùa Xuân đến, cả bốn người bọn tôi có thể đứng cạnh nhau một lần nữa.

\il

Tôi tựa người vào ghế, nhìn sang hai đồng đội của mình - những người cũng tới trú mưa, chờ đợi mưa tạnh. Bây giờ, khi mọi chuyện đã tạm ổn thỏa, cả ba có thể ngồi lại và tổng kết lại những gì đã diễn ra trong ngày đầu tiên chúng tôi đặt chân lên văn phòng này.

``Cảm giác hôm nay\dots\ đúng là hỗn loạn thật sự,'' tôi lên tiếng trước, vừa nói vừa xoa thái dương.

Polka cười khúc khích: ``Công nhận Shishiron nói đúng. Bọn mình mới vào công ty thôi mà đã có một đống chuyện điên rồ xảy ra rồi.''

Lamy gật gù, đưa tay vuốt ve Daifuku, như để trấn an chính mình.

Những tuần vừa rồi đúng là một chuỗi sự kiện điên rồ.

Polka đã bị Fubuki-senpai cùng Haato-senpai và Ayame-senpai thuyết phục rằng cậu ta là một con mèo. Rõ ràng, tôi còn nhớ là nhóm tôi có cáo, có sư tử, nhưng chẳng có ai là mèo cả. Nhưng cuối cùng, sự trêu đùa của ba người đó đã khiến Fubuki-senpai bị Okayu-senpai và Korone-senpai bắt cóc vì tưởng nhầm chị ấy cũng là mèo.

Kết quả là Polka phải hợp sức với ba người để giải cứu chị ấy. Đến cuối cùng, chuyện còn đi xa hơn nữa khi mọi người quyết định biến nó thành một trò chơi nhập vai thực sự, với Okayu-senpai và Korone-senpai làm nhân vật chính, do chính Kuon-senpai bày trò.

Còn Lamy\dots\ Ôi trời, chuyện của cậu ấy đúng là một bộ phim hài riêng biệt.

Cậu ấy bị Mio-sepai, Flare-senpai và Kuon-sepai ``tra khảo'' về số rượu mà cậu ấy mang vào. Đúng là cậu ấy nghiện rượu nặng thật. Và rồi, trong một nỗ lực tuyệt vọng để thoát khỏi tình huống khó xử, cậu ấy đã triệu hồi Daifuku. Chỉ có điều\dots\ đồng đội lại phản bội cô ấy ngay lập tức bằng cách giật mất chai rượu quý kia. Lamy khóc lóc cầu cứu, nhưng không ai đứng về phía cậu ấy cả.

Chưa hết, anh Tổng của chúng ta còn trêu thêm bằng cách rủ chính đồng đội của cậu ấy đi nhậu tại nhà mình. Nhìn cảnh Lamy tức giận nhưng chẳng thể làm gì đúng là buồn cười chết đi được. Nhưng sau một hồi giằng co, cuối cùng anh ấy cũng đồng ý cho Lamy đi theo.

Và thế là cậu ấy có một bữa nhậu tưng bừng với gia đình Kuon-senpai, thậm chí còn dám thách đấu tửu lượng với bố của anh ấy. Đúng là cái gan của ``công chúa tuyết'' này cũng không phải dạng vừa đâu.

Còn tôi\dots\ Tôi thì được Kuon hòa giải giữa tôi, Towa-senpai và Luna về cái mà em ấy gọi là ``triết lý xe ba bánh''. Nguyên nhân chỉ đơn giản là do một khẩu súng mô hình mà tôi đã lắp ráp, vậy mà rốt cuộc lại thành một trận rượt đuổi trong văn phòng. Nghĩ lại, tôi vẫn cảm thấy nhức đầu.

Buổi họp bất thường  kết thúc trong một bầu không khí lặng lẽ hơn hẳn so với mọi khi. Trên màn hình điện thoại, khuôn mặt quen thuộc của người đồng đội thứ tư hiện lên, nhưng không phải với nụ cười rạng rỡ như thường thấy. Trông cậu ta vẫn còn khá xanh xao.

``Xin lỗi mọi người\dots\ Hôm nay mình vẫn không lên văn phòng được\dots'' Giọng cậu ấy yếu ớt, có chút mệt mỏi.

``Cậu vẫn chưa khỏi hẳn à?'' Lamy lo lắng hỏi.

Cô ấy lắc đầu nhẹ: ``Không hẳn\dots\ Mình thì đỡ rồi, nhưng cũng không hẳn chỉ vì bệnh đâu\dots''

Không ai trong bọn tôi lên tiếng. Dù cô ấy không nói thẳng, nhưng cả ba chúng tôi đều hiểu.

Vì sao cậu ấy lại bệnh?

Vì sao cậu ấy lại không thể lên văn phòng?

Tất cả không chỉ là do thời tiết.

``Dù đã bao lâu trôi qua, tớ vẫn không thể quên được\dots'' cậu ấy nói, giọng trầm xuống. ``Người đó\dots\ người đã từng sát cánh bên chúng ta\dots''

Căn phòng chợt trở nên yên ắng một cách kỳ lạ.

Không ai muốn nói ra, nhưng chúng tôi đều biết cậu ấy đang nhắc đến ai.

Một mảnh ghép từng thuộc về \textit{holofive}, từng là một phần của nhóm chúng tôi, giờ đã không còn ở đây nữa.

Lamy cúi đầu, những ngón tay vô thức siết chặt vạt áo. Polka im lặng, đôi mắt nhìn xuống nền nhà, như thể đang cố giấu đi điều gì đó. Còn tôi\dots\ cảm thấy một thứ gì đó nghẹn lại trong lòng, nhưng lại chẳng thể thốt nên lời.

Mỗi người chúng tôi đều tỏ vẻ một cảm xúc riêng, nhưng, có một điểm chung, đó là chúng tôi sẽ mãi người bạn đó. Sau tai nạn đáng tiếc ngày \emph{hôm đó}, chúng tôi đã không còn yếu đuối như ngày trước được nữa.

Chúng tôi giờ đã trở thành một tập thể bốn người với sự gắn kết mạnh mẽ. Chúng tôi đã trở nên khăng khít hơn bao giờ hết. Tôi - với tư cách là nhóm trưởng (thực tế) của \textit{NePoLaBo} đây - sẽ không để tai nạn đó lặp lại với bất cứ ai trong chúng tôi thêm một lần nào nữa.

Sau một khoảng thời gian im lặng chìm vào những nỗi suy tư của chúng tôi, cuối cùng, Lamy khẽ cười, một nụ cười không còn sự vui vẻ như thường ngày. Cậu ấy bất giác nói: ``Một ngày nào đó\dots\ năm người chúng ta sẽ lại tái ngộ nhau chứ?''

Polka ngẩng lên, đôi mắt cậu ấy phản chiếu ánh đèn từ màn hình điện thoại. ``Ừm\dots\ dù cho điều đó có bất khả thi đến đâu.'' Giọng của Omarun nhẹ bẫng, gần như một tiếng thở dài.

Tôi chỉ lặng lẽ gật đầu.

Tôi không biết chuyện đó có thể thành sự thật hay không, nhưng tôi vẫn mong rằng, nếu có một ngày như thế\dots\ tất cả chúng tôi sẽ lại được đứng cạnh nhau, như những ngày đầu tiên.

\pov{Hikawa Kuon}

Hôm nay đáng lẽ đã có thể kết thúc một cách bình yên.

Sau khi lo xong mọi chuyện, từ giải quyết rắc rối của Botan-chan, Luna-chan và Towa-chan, đến cả chuyện của Marine-chan, tôi cuối cùng cũng có thể rời công ty để về nhà. Trời đã tối, nhưng ánh đèn đường vẫn sáng rực, phản chiếu lên nền đất ướt sau cơn mưa chiều nay.

Khổ thật chứ, tới lúc tôi về rồi, trời vẫn còn mưa trắng trời, hầu như không có dấu hiệu dừng lại. Tôi lại còn không mang ô theo nữa. Đúng là tôi biết rằng trời hôm nay sẽ mưa, nhưng đâu đến mức dai dẳng như thế này!

``Khỉ thật!'' tôi chẹp miệng, dùng luôn cái cặp máy tính chạy một mạch từ tòa nhà văn phòng tới trạm xe buýt. Ơn giời là nó rất gần với văn phòng, có mái che, và cặp tôi cũng là cặp chống nước, chứ nếu không thì tôi cũng chôn chân ở đó, hoặc có nếu liều thì một là hỏng toi cái máy, hai là cả người tôi sẽ ướt như chuột lội\footnote{Thành ngữ gốc và chuẩn xác nhất của ``ướt như chuột lột.''} mất.

Như tôi đã nói, hôm nay là ngày làm việc cuối cùng của tôi trước khi nghỉ Tết Nguyên Đán. Giờ, tôi chỉ cần về nhà, chuẩn bị mâm cỗ Tết, dọn dẹp nhà cửa, và sẽ lại đón thêm một kỳ nghỉ lễ nữa. Đây sẽ là kỳ nghỉ Tết thứ hai của tôi kể từ khi tôi nhận vào \hll\ này.

Tôi ngồi đó, định lấy điện thoại ra đọc báo giết thời gian chờ chuyến xe của tôi về nhà, thì bỗng ông sếp YAGOO - hay Tanigo-san - gọi cho tôi. Tôi chắc mẩm: ``\textit{Có khi ông ấy định hỏi thăm với chúc Tết mình sớm đây.}''

Chẳng ngờ rằng, cuộc gọi đó trở thành cuộc gọi định mệnh không chỉ với tôi, mà cả gia đình Hikawa chúng tôi. Sau cuộc gọi đó, tôi vẫn còn bị choáng, không thể tin được rằng\dots

\dots gia đình Hikawa chúng tôi sẽ sang một ngã rẽ hoàn toàn khác. Sự hỗn loạn mà tôi chịu đựng và đối mặt hàng ngày, hàng tuần, hàng tháng trời tưởng chừng sẽ chỉ dừng lại ở thành phố Điện tử này, nhưng có vẻ như nó sẽ còn được lan sang cả Kawachi - ngay tại chính quê hương của tôi, ngay tại chính nhà Hikawa chúng tôi.

Nhưng, chờ một chút. Đây lại là một câu chuyện khác mà tôi sẽ kể trong **Quyển D: Ký túc xá nhà Hikawa.**

\end{document}